\section{Conclusion}
\label{sec:conclusion}

We presented an offline system that turns long-shot windsurfing video into stable,
rider-relative trajectories.  Its design follows the structure of the problem:
masked camera-motion estimation handles large pans, conservative local association
protects tracklet purity, sail-specific appearance and four-dimensional motion score
longer continuations, a global integer program uses complete-video context, and two
rig landmarks control the final virtual camera.

On 21 manually reconstructed development videos, the association subsystem improves
pairwise F1 and greatly reduces fragmentation relative to the evaluated OC-SORT and
BoT-SORT configurations.  The staged revision also shows why implementation details
matter: activating the intended foreground mask, aggregating appearance over whole
fragments, retaining box-size evidence, and vetoing implausible links collectively
reduce contaminated tracks from 20 to nine.  The remaining crossing errors and the
in-sample protocol preclude a claim of perfect or general-purpose tracking.  The
contribution is instead a reproducible account of how domain structure, offline
inference, and an explicit failure cost can make an otherwise ordinary detection
model useful in a demanding video-analysis system.
